\section{Architecture}

\subsection{Who this guide is for}

This guide assumes you have never opened the \spetc{} source. It explains
what the program is made of, how a calculation flows through it, where
every kind of file lives, and then --- once the map is in your head ---
how to change things safely: fix a bug, add a physical effect, add data,
publish a release. Read this section and Sect.~\ref{sec:modules} before
touching anything; the recipes of Sect.~\ref{sec:extending} then tell you
exactly which file to open for each kind of change.

\subsection{The shape of the program}

\spetc{} is deliberately \emph{not} a framework: it is sixteen flat Python
modules in one directory, no packages, no plugins, no configuration
language. Two classes of module exist, distinguished by one strict rule:

\begin{quote}
\textbf{Physics modules never import the GUI, and only \code{etc\_gui.py}
imports Tkinter.}
\end{quote}

The consequence is that everything scientific is importable, testable and
scriptable from a plain Python session or the test suite, with no window
involved; the GUI is a thin, replaceable skin. The batch tool
\code{spetc\_batch.py} is the proof: it produces complete results using
the same engines and zero Tk.

Within the physics side, the modules form five layers, each depending
only on the layers below it:

\begin{enumerate}
\item \textbf{Foundations.} \code{spectral\_utils.py} defines the one
  data structure the whole program shares: a \code{SpectralCurve}, a
  validated pair of (strictly increasing wavelength in \AA{}, value)
  arrays, with three interpolation policies that encode the project's
  coverage philosophy --- \code{interpolate\_checked} (raise outside
  coverage: for calibration), \code{interpolate\_zero\_filled} (assume
  zero outside: for detector-side quantities like QE),
  \code{interpolate\_edge\_extended} (extend the edge value: for the
  atmosphere, which does not become opaque where a file ends). It also
  reads two-column text files and atmospheric FITS tables.
  \code{observing\_conditions.py} is a set of pure functions for the
  atmosphere-as-noise: scintillation (Young's law, in both per-exposure
  and per-second-of-variance form), Filippenko differential refraction,
  and ADC quantization noise. Neither imports anything else of the
  project except each other.
\item \textbf{Physics core.} \code{etc\_physics.py} is the radiometric
  heart: magnitude systems and zero points, synthetic photometry in the
  SVO convention, the template flux calibration (the $m_{v0}$ machinery),
  the electron-rate integral, the PSF models (Gaussian and Moffat, with
  their aperture and slit couplings) and the CCD signal-to-noise
  equation. \code{sky\_brightness.py} and \code{sky\_background.py}
  together model the sky: the first owns the component physics (van
  Rhijn airglow, zodiacal light, starlight, the Krisciunas--Schaefer
  Moon) and the shared constant tables; the second owns the
  \emph{normalization policy} --- which component sets the level at a
  given Sun altitude (dark table, SQM reading, Weaver daylight, twilight
  blend). \code{ephemeris.py} wraps Astropy for everything positional:
  refraction-corrected altitude--azimuth tracks, Pickering airmass,
  Sun/Moon positions and phase, parallactic angle, and coordinate/time
  parsing.
\item \textbf{Engines.} \code{photometry.py} and \code{spectroscopy.py}
  each expose one class with one compute method. They accept what we call
  the \emph{four-dictionary context} --- \code{telescope} (a plain dict
  of numbers and optional curves), \code{detector} (a \code{Detector}
  instance from \code{detector.py}), \code{atmosphere} (airmass, seeing,
  PSF model, elevation, transmission curve), \code{sky\_model}
  (brightness, calibration, geometry) --- plus per-call arguments, and
  return plain results: a dict for photometry, a \code{pandas.DataFrame}
  with metadata in \code{.attrs} for spectroscopy. This context is the
  API contract of the whole program: anything that can build those four
  objects can run the ETC. \code{solvers.py} inverts the S/N equation in
  closed form and plans frame stacks.
\item \textbf{Catalogues.} \code{star\_catalog.py} parses the template
  catalogue and loads spectra (ASCII and CALSPEC-style FITS);
  \code{filter\_catalog.py} parses SVO VOTable filter profiles and lists
  the available filters; \code{config\_manager.py} loads the built-in
  observatory presets.
\item \textbf{Presentation and tools.} \code{etc\_gui.py} (the Tk
  application) and \code{matplotlib\_plots.py} (the native plot windows)
  on the interactive side; \code{spetc\_batch.py},
  \code{make\_filter\_profile.py} and \code{add\_template.py} as
  command-line tools; \code{self\_check.py} and
  \code{tests/test\_spectral\_pipeline.py} as the test entry points.
\end{enumerate}

\subsection{One calculation, end to end}

To understand the code, follow one press of \textsf{Run ETC} through it.

\begin{enumerate}
\item \textbf{Input validation.} \code{ETCGUI.\_run\_etc} reads every
  relevant Tk variable, converts and range-checks it, and constructs the
  target coordinates. Anything invalid raises \code{ValueError} with a
  human sentence; the GUI catches it and shows it --- this is the
  program's error policy everywhere: \emph{raise early, with the
  offending number in the message, and never substitute a silent
  default}.
\item \textbf{Ephemerides.} \code{ephemeris.compute\_target\_track}
  produces, for a window of $\pm$18 hours around the chosen time, the
  refraction-corrected altitude, azimuth, Pickering airmass, parallactic
  angle, and the Sun and Moon positions, separations and phase at
  five-minute steps. Everything time-dependent downstream is evaluated on
  this track.
\item \textbf{Sky models.} \code{\_sky\_models\_for\_track} converts the
  target's ICRS coordinates to ecliptic and galactic latitude once, then
  builds one sky-model dictionary per time sample: the observed surface
  brightness in the observing band (from the ING physical model, the
  user's SQM reading, or a fixed value, plus the Moon term), the
  photometric zero point it is expressed in, and the spectral colour
  model used by spectroscopy.
\item \textbf{Engine loop.} For every visible time sample, the GUI builds
  the four-dictionary context and calls the engine. In target-S/N mode it
  first calls the engine at 1\,s to obtain rates, inverts the S/N
  equation in closed form (\code{solvers.exposure\_time\_for\_snr},
  scintillation included as an extra variance rate), and calls the
  engine again at the solved exposure. Full spectra are computed only at
  the selected time and at the plot-slider samples; all other samples
  evaluate just the reference-wavelength bin, which keeps a night's
  planning at $R = 100\,000$ responsive.
\item \textbf{Aggregation and display.} The per-time values become the
  time-series DataFrame (one row per sample: time, geometry, parallactic
  angle, airmass, S/N or required exposure, sky brightness, extinction,
  saturation); the selected-time result becomes the results table; the
  stack planner and the differential-precision helper run on the
  selected-time rates; plots and the parameters/outputs log are
  refreshed; the configuration is saved.
\end{enumerate}

\subsection{Inputs and outputs at a glance}

\begin{center}
\begin{tabular}{llp{6.6cm}}
\toprule
File & Direction & Role \\
\midrule
\code{data/interpola.db.csv} & in & template catalogue (GUI list) \\
\code{data/bpgs/*}, \code{data/imported/*} & in & template spectra (ASCII/FITS) \\
\code{data/filters.list}, \code{data/filters/*.xml} & in & filter profiles (SVO VOTable) \\
\code{data/qe.dat}, curve files & in & QE, optics response, slit-$R$, gain table \\
\code{data/earth\_atmospheric\_transmission.fits} & in & measured zenith transmission \\
\code{etc\_sites.json}, \code{observatories.json} & in/out & site lists \\
\code{etc\_user\_config.json} & in/out & full GUI state, rewritten each run \\
instrument profile JSON + curve copies & in/out & portable per-instrument setup \\
result CSV / time-series CSV & out & full numerical results (every engine field) \\
batch JSON & in & headless run description \\
batch result CSV + HTML summary & out & headless outputs \\
\code{dem\_cache/*.tif} & in/out & cached Copernicus DEM tiles \\
\code{data/horizons/*.csv} & out & generated horizon profiles \\
\bottomrule
\end{tabular}
\end{center}

The complete field-by-field specification of every input format is in the
\emph{User Guide}, Sect.~6, written to be sufficient for producing the
files from scratch; this guide does not duplicate it, but
Sect.~\ref{sec:formats} adds the parsing internals a maintainer needs.

\subsection{Project conventions}

Wavelengths are \AA{}ngstr\"om internally, everywhere; every user-facing
curve declares its unit explicitly and is converted once at load. All
radiometric integrals carry \code{astropy.units} and are reduced to
floats only at the engine boundary. Constants are named after their
paper (\code{SOLAR\_MINIMA}, \code{MOON\_FUDGE}); every empirical
relation cites its reference in the docstring; every physical claim in
the code has a regression test pinned to a published number, not to the
code's own output. Azimuth is N$=0^\circ$/E$=90^\circ$; times are UTC
internally and converted only for display.
